\documentclass[12pt,aspectratio=169]{beamer}
\input{header_beamer}
\usetheme{Boadilla}
\usecolortheme{beaver}
\setbeamertemplate{page number in head/foot}{}
\setbeamertemplate{navigation symbols}{}
\title{Lecture-13: Conditional expectation (Continued)}
\date{}

\begin{document}
\pagestyle{empty}
\maketitle


\begin{frame}{1 Conditional expectation conditioned on a random  vector}
\begin{itemize}
\item <1-> Let $X: \Omega \to \R^n$ be a random vector defined on a probability space $(\Omega, \sF, P)$. 
\item <2-> Recall that the projection $\pi_i: \R^n \to \R$ of $n$-dimensional vector $X$ on its $i$th component gives $X_i = \pi_i(X)$, 
and it is a Borel measurable function. 
It follows that $X_i = \pi_i\circ X$ is a random variable for each $i \in [n]$. 
\item <3-> Further the event space generated by random vectors $X$ is given by 
\EQ{
\sigma(X) = \sigma(X_1, \dots, X_n).
}
\end{itemize}
\begin{block}{Reminder}<4->
Let $(A_i \in \sF:  i \in [n])$ be an $n$-length sequence of events, 
then $X = (\Ind{A_i}: i \in [n])$ is a random vector, and the smallest event space generated by this random vector is $\sigma(X) = \sigma(A_1, \dots, A_n)$. 
\end{block}
\end{frame}

\begin{frame}
\begin{Definition}[Conditional expectation given a random vector]<1->
Consider a random variable $X:\Omega \to\R$ such that $\E\abs{X} < \infty$ and a random vector $Y:\Omega \to \R^n$ for some $n\in \N$, 
defined on the same probability space $(\Omega, \sF, P)$. 
The \textbf{conditional expectation of random variable $X$ given the random vector $Y$} is defined as $\E[X\mid Y] \triangleq \E[X\mid\sigma(Y)]$. 
\end{Definition}
\begin{block}{Lemma}<2->
For a random sequence $X: \Omega \to \R^\N$ defined on the same probability space $(\Omega, \sF, P)$, 
we have 
\EQ{
\sigma(X_1, \dots, X_n) \subseteq \sigma(X_1, \dots, X_{n+1}),\quad n\in \N.
}
\end{block}

\begin{Proof}<3->
For any $x \in \R^{n+1}$, 
any generating event for collection $\sigma(X_1, \dots, X_n)$ is of the form $\cap_{i=1}^nX_i^{-1}(-\infty, x_i] = \cap_{i=1}^nX_i^{-1}(-\infty, x_i]\cap X_{n+1}^{-1}(\R)$, a generating event for collection $\sigma(X_1, \dots, X_{n+1})$.  
\end{Proof}
\end{frame}


\begin{frame}{2 Properties of Conditional Expectation}
\begin{prop} 
Let $X,Y$ be random variables on the probability space $(\Omega, \sF, P)$ such that $\E\abs{X} , \E\abs{Y} < \infty$. 
Let $\sG$ and $\sH$ be event spaces such that $\sG, \sH \subset \sF$. 
Then the following statements are true. 
\begin{enumerate}
\item <2-> \textbf{Identity:} If $X$ is $\sG$-measurable and $\E\abs{X} < \infty$, then $X = \E[X\mid\sG]$ a.s. 
In particular, $c = \E[c\mid \sG]$, for any constant $c \in  \R$. 
\end{enumerate}
\end{prop}
\end{frame}


\begin{frame}
\begin{Proof}
\begin{enumerate}
\item <1-> \textbf{Identity:} It follows from the definition that $X$ satisfies all three conditions for conditional expectation.  
The event space generated by any constant function is the trivial event space $\set{\emptyset, \Omega} \subseteq \sG$ for any event space. 
Hence, $\E[c\mid \sG] = c$. 
\end{enumerate}
\end{Proof}
\end{frame}

\begin{frame}{2 Properties of Conditional Expectation}
\begin{prop} 
Let $X,Y$ be random variables on the probability space $(\Omega, \sF, P)$ such that $\E\abs{X} , \E\abs{Y} < \infty$. 
Let $\sG$ and $\sH$ be event spaces such that $\sG, \sH \subset \sF$. 
Then the following statements are true. 
\begin{enumerate}
\item \textbf{Identity:} If $X$ is $\sG$-measurable and $\E\abs{X} < \infty$, then $X = \E[X\mid\sG]$ a.s. 
In particular, $c = \E[c\mid \sG]$, for any constant $c \in  \R$. 
\item \textbf{Linearity:} $\E[(\alpha X + \beta Y) \mid \sG] = \alpha \E[X\mid\sG] + \beta E[Y\mid \sG]$ a.s. 
\end{enumerate}
\end{prop}
\end{frame}


\begin{frame}
\begin{Proof}
\begin{enumerate}
\setcounter{enumi}{1}
\item \textbf{Linearity:} 
Let $Z \triangleq \E[(\alpha X + \beta Y) \mid \sG]$ and $Z_1 \triangleq \E[X\mid\sG]$ and $Z_2\triangleq\E[Y\mid\sG]$. 
We have to show that $Z = \alpha Z_1+ \beta Z_2$ almost surely. 
\begin{enumerate}[(i)]
\item <2-> From the definition of conditional expectation, we have $\E[X\mid\sG], \E[Y\mid\sG]$ are $\sG$-measurable. 
In addition, the linear map $g: \R^2\to\R$ defined by $g(x) \triangleq \alpha x_1 + \beta x_2$ is Borel measurable. 
Therefore, the linear combination $\alpha Z_1 + \beta Z_2$ is $\sG$-measurable.
\item <3-> For any $G \in \sG$, from the linearity of expectation and definition of conditional expectation, we have 
\EQ{
\E[(\alpha Z_1 + \beta Z_2)\Ind{G}] 
= \alpha\E[ \E[X\mid \sG]\Ind{G}] + \beta\E[ \E[Y\mid \sG]\Ind{G}] 
= \E[(\alpha X + \beta Y)\Ind{G}].
} 
\item <4-> From the definition of conditional expectation, we have $\E\abs{Z_1}, \E\abs{Z_2}$ are finite. 
Therefore, we have 
\EQ{
\E\abs{\alpha Z_1 + \beta Z_2} \le \abs{\alpha}\E\abs{Z_1} + \abs{\beta}\E\abs{Z_2} < \infty.
}
\end{enumerate}
%This implies that $\alpha Z_1 + \beta Z_2$ satisfies three properties of conditional expectation $\E[(\alpha X +\beta Y)\mid\sG]$. 
%From the almost sure uniqueness of conditional expectation, we have $Z = \alpha Z_1 + \beta Z_2$ almost surely. 
\end{enumerate}
\end{Proof}
\end{frame}


\begin{frame}{2 Properties of Conditional Expectation}
\begin{prop} 
Let $X,Y$ be random variables on the probability space $(\Omega, \sF, P)$ such that $\E\abs{X} , \E\abs{Y} < \infty$. 
Let $\sG$ and $\sH$ be event spaces such that $\sG, \sH \subset \sF$. 
Then the following statements are true. 
\begin{enumerate}
\item \textbf{Identity:} If $X$ is $\sG$-measurable and $\E\abs{X} < \infty$, then $X = \E[X\mid\sG]$ a.s. 
In particular, $c = \E[c\mid \sG]$, for any constant $c \in  \R$. 
\item \textbf{Linearity:} $\E[(\alpha X + \beta Y) \mid \sG] = \alpha \E[X\mid\sG] + \beta E[Y\mid \sG]$ a.s.
\item \textbf{Monotonicity:} If $X \le Y$ a.s., then $\E[X\mid\sG] \le E[Y\mid \sG]$ a.s.
\end{enumerate}
\end{prop}
\end{frame}


\begin{frame}
\begin{Proof}[Proof Continued.]
\begin{enumerate}
\setcounter{enumi}{2}
\item \textbf{Monotonicity:} Let $\epsilon > 0$ and define $A_\epsilon \triangleq \set{\E[X\mid\sG] - \E[Y\mid \sG] >\epsilon} \in  \sG$.  
Then from the definition of conditional expectation, we have 
\EQ{
0 \le \epsilon P(A_\epsilon) 
= \E[\epsilon\Ind{A_\epsilon}] 
< \E[(\E[X\mid \sG]-\E[Y\mid\sG])\Ind{A_\epsilon}] 
= \E[(X-Y)\Ind{A_\epsilon}] 
\le 0.
}
Thus, we obtain  that $P(A_\epsilon) = 0$ for all $\epsilon > 0$. 
Taking limit $\epsilon \downarrow 0$, we get 
\EQ{
0 =\lim_{\epsilon \downarrow 0}P(A_\epsilon) = P(\lim_\epsilon A_\epsilon) = P(A_0).
} 
\end{enumerate}
\end{Proof}
\end{frame}

\begin{frame}{2 Properties of Conditional Expectation}
\begin{prop} 
Let $X,Y$ be random variables on the probability space $(\Omega, \sF, P)$ such that $\E\abs{X} , \E\abs{Y} < \infty$. 
Let $\sG$ and $\sH$ be event spaces such that $\sG, \sH \subset \sF$. 
Then the following statements are true. 
\begin{enumerate}
\item \textbf{Identity:} If $X$ is $\sG$-measurable and $\E\abs{X} < \infty$, then $X = \E[X\mid\sG]$ a.s. 
In particular, $c = \E[c\mid \sG]$, for any constant $c \in  \R$. 
\item \textbf{Linearity:} $\E[(\alpha X + \beta Y) \mid \sG] = \alpha \E[X\mid\sG] + \beta E[Y\mid \sG]$ a.s.
\item \textbf{Monotonicity:} If $X \le Y$ a.s., then $\E[X\mid\sG] \le E[Y\mid \sG]$ a.s.
\item \textbf{Conditional Jensen's inequality:} If $\psi : \R \to  \R$ is convex and $\E\abs{\psi(X)} <\infty$, 
then $\psi(\E[X\mid\sG]) \le \E[\psi(X) \mid \sG]$ a.s. 
\end{enumerate}
\end{prop}
\end{frame}


\begin{frame}
\begin{Proof}[Proof Continued.]
\begin{enumerate}
\setcounter{enumi}{3}
\item \textbf{Conditional Jensen's inequality:}
We will use the fact that a convex function can always be represented by the supremum of a family of affine functions. 
Accordingly, we will assume for a convex function $\psi:\R\to\R$, 
we have linear functions $\phi_i:\R\to\R$ and constants $c_i\in\R$ for all $i \in I$ such that $\psi = \sup_{i\in I}(\phi_i+c_i)$. 

For each $i\in I$, we have $\phi_i(\E[X\mid\sG])+c_i = \E[\phi_i(X)\mid\sG] +c_i\le \E[\psi(X)\mid\sG]$ from the linearity and monotonicity of conditional expectation. 
It follows that
\EQ{
\psi(\E[X\mid\sG]) = \sup_{i \in I}(\phi_i(\E[X\mid\sG]) + c_i) \le \E[\psi(X)\mid\sG].
}
\end{enumerate}
\end{Proof}
\end{frame}


\begin{frame}
\begin{prop} [Properties continued.]
Let $X,Y$ be random variables on the probability space $(\Omega, \sF, P)$ such that $\E\abs{X} , \E\abs{Y} < \infty$. 
Let $\sG$ and $\sH$ be event spaces such that $\sG, \sH \subset \sF$. 
Then the following statements are true. 
\begin{enumerate}
\setcounter{enumi}{4}
\item \textbf{Pulling out what's known:} If $Y$ is $\sG$-measurable and $\E\abs{XY} <\infty$, then
$\E[XY\mid \sG] = Y\E[X\mid\sG]$ a.s.
\end{enumerate}
\end{prop}
\end{frame}


\begin{frame}
\begin{Proof}[Proof Continued.]
\begin{enumerate}
\setcounter{enumi}{4}
\item \textbf{Pulling out what's known:} 
From the almost sure uniqueness of conditional expectation,  
it suffices to show that $Y\E[X\mid\sG]$ satisfies following three properties of the conditional expectation $\E[XY\mid\sG]$: 
\begin{enumerate}[(i)]
\item the random variable $Y\E[X\mid\sG]$ is $\sG$-measurable,
\item $\E[XY\Ind{G}] = \E[Y\E[X\mid\sG]\Ind{G}]$ for any event $G\in\sG$, and 
\item $\E\abs{Y\E[X\mid\sG]}$ is finite.
\end{enumerate}
Part~(i) is true since $Y$ is given to be $\sG$-measurable, 
$\E[X\mid\sG]$ is $\sG$-measurable by the definition of conditional expectation, 
and product function is Borel measurable.
\end{enumerate}
\end{Proof}
\end{frame}


\begin{frame}
\begin{Proof}[Proof Continued.]
\begin{enumerate}
\setcounter{enumi}{4}
\item It suffices to show part~(ii) and~(iii) for simple $\sG$-measurable random variables $Y=\sum_{y\in\sY}y\Ind{E_y}$ where $E_y = Y^{-1}\set{y} \in \sG$. %, where $\E\abs{XY}$ is finite. 
For any $G\in\sG$, we have $G\cap E_y\in\sG$ for all $y \in \sY$ and $\cup_{y\in\sY}(G\cap E_y) = G$. 

Part~(ii) follow from the linearity of expectation and definition of conditional expectation $\E[X\mid\sG]$, 
such that  
\EQ{
\E[Y\E[X\mid\sG]\Ind{G}] 
= \sum_{y\in\sY}y\E[\Ind{G\cap E_y}\E[X\mid\sG]]
= \sum_{y\in\sY}y\E[X\Ind{G\cap E_y}]
= \E[X\sum_{y\in\sY}y\Ind{G\cap E_y}]
= \E[XY\Ind{G}].
}

Part~(iii) follows from the fact that $\E\abs{XY}$ is finite and the conditional Jensen's inequality applied to convex function $\abs{\cdot} : \R \to \R_+$ to get $\abs{\E[X\mid\sG]} \le \E[\abs{X}\mid\sG]$. 
Therefore, 
\EQ{
\E[\abs{Y}\abs{\E[X\mid\sG]}] 
= \sum_{y\in\sY}\abs{y}\E[\abs{\E[X\mid\sG]}\Ind{E_y}]
\le \sum_{y \in \sY}\abs{y}\E[\abs{X}\Ind{E_y}] 
= \E\abs{XY}.  
} 
\end{enumerate}
\end{Proof}
\end{frame}


\begin{frame}
\begin{prop} [Properties continued.]
Let $X,Y$ be random variables on the probability space $(\Omega, \sF, P)$ such that $\E\abs{X} , \E\abs{Y} < \infty$. 
Let $\sG$ and $\sH$ be event spaces such that $\sG, \sH \subset \sF$. 
Then the following statements are true. 
\begin{enumerate}
\setcounter{enumi}{4}
\item \textbf{Pulling out what's known:} If $Y$ is $\sG$-measurable and $\E\abs{XY} <\infty$, then
$\E[XY\mid \sG] = Y\E[X\mid\sG]$ a.s.
\item \textbf{Tower property:} If $\sH \subseteq \sG$, then
$\E[\E[X\mid\sG]\mid \sH] = \E[X\mid \sH]$ a.s.
\end{enumerate}
\end{prop}
\end{frame}


\begin{frame}
\begin{Proof}[Proof Continued.]
\begin{enumerate}
\setcounter{enumi}{5}
\item \textbf{Tower property:} 
From the almost sure uniqueness of conditional expectation, 
it suffices to show that 
\begin{enumerate}[(i)]
\item $\E[X\mid\sH]$ is $\sH$-measurable, 
\item $\E[\E[X\mid\sH]\Ind{H}] = \E[\E[X\mid\sG]\Ind{H}]$ for all $H\in\sH$, and 
\item $\E[\abs{\E[X\mid\sH]}]$ is finite.  
\end{enumerate}
Part~(i) follows from the definition of conditional expectation, 
which implies that $\E[X\mid \sH]$ is $\sH$ measurable. 

Part~(ii) follows from the fact that $\sH \subseteq \sG$, 
and hence any $H \in \sH$ belongs to $\sG$. 
Therefore, from the definition of conditional expectation, we have 
\EQ{
\E[\E[X\mid \sG]\Ind{H}] = \E[X\Ind{H}] = \E[\E[X\mid \sH]\Ind{H}]. 
}

Part~(iii) follows from the conditional Jensen's inequality applied to convex function $\abs{\cdot}:\R\to\R_+$ to get $\abs{\E[X\mid\sH]} \le \E[\abs{X}\mid\sH]$, 
and hence $\E\abs{\E[X\mid\sH]} \le \E\abs{X} < \infty$. 
\end{enumerate}
\end{Proof}
\end{frame}


\begin{frame}
\begin{prop} [Properties continued.]
Let $X,Y$ be random variables on the probability space $(\Omega, \sF, P)$ such that $\E\abs{X} , \E\abs{Y} < \infty$. 
Let $\sG$ and $\sH$ be event spaces such that $\sG, \sH \subset \sF$. 
Then the following statements are true. 
\begin{enumerate}
\setcounter{enumi}{4}
\item \textbf{Pulling out what's known:} If $Y$ is $\sG$-measurable and $\E\abs{XY} <\infty$, then
$\E[XY\mid \sG] = Y\E[X\mid\sG]$ a.s.
\item \textbf{Tower property:} If $\sH \subseteq \sG$, then
$\E[\E[X\mid\sG]\mid \sH] = \E[X\mid \sH]$ a.s.
\item \textbf{Irrelevance of independent information:} If $\sH$ is independent of $\sigma (\sG, \sigma (X))$
then
$\E[X|\sigma (\sG, \sH)] = \E[X\mid\sG] \text{ a.s. }$
In particular, if $X$ is independent of $\sH$, then $\E[X\mid \sH] = \E[X]$ a.s.
\end{enumerate}
\end{prop}
\end{frame}


\begin{frame}
\begin{Proof}[Proof Continued.]
\begin{enumerate}
\setcounter{enumi}{6}
\item \textbf{Irrelevance of independent information:} 
From the almost sure uniqueness of conditional expectation, 
it suffices to show that 
\begin{enumerate}[(i)]
\item $\E[X\mid\sG]$ is $\sigma(\sG,\sH)$-measurable, 
\item $\E[\E[X\mid\sG]\Ind{G}] = \E[X\Ind{G}]$ for all $G\in\sigma(\sG,\sH)$, and 
\item $\E\abs{\E[X\mid\sG]}$ is finite.  
\end{enumerate}

Part~(i) Since $\E[X\mid\sG]$ is $\sG$-measurable, it is $\sigma(\sG,\sH)$ measurable. 


Part~(ii) For events $A = G \cap H \in \sigma(\sG,\sH)$ 
where $G \in \sG$ and $H \in \sH$,  
\begin{align*}  
\E[\E[X\mid\sG]\Ind{G\cap H}] 
& = \E[\E[X\mid\sG]\Ind{G}\Ind{H}] 
= \E[\E[X\mid\sG]\Ind{G}]E[\Ind{H}] \\
& = \E[X\Ind{G}]\E[\Ind{H}] = \E[X\Ind{G\cap H}]. 
\end{align*}

Part~(iii) From the conditional Jensen's inequality applied to $\abs{\E[X\mid\sG]} \le \E[\abs{X}\mid\sG]$,  
\EQ{\E\abs{\E[X\mid\sG]} \le \E\abs{X} < \infty}. 
\end{enumerate}
\end{Proof}
\end{frame}


\begin{frame}
\begin{prop} [Properties continued.]
Let $X,Y$ be random variables on the probability space $(\Omega, \sF, P)$ such that $\E\abs{X} , \E\abs{Y} < \infty$. 
Let $\sG$ and $\sH$ be event spaces such that $\sG, \sH \subset \sF$. 
Then the following statements are true. 
\begin{enumerate}
\setcounter{enumi}{4}
\item \textbf{Pulling out what's known:} If $Y$ is $\sG$-measurable and $\E\abs{XY} <\infty$, then
$\E[XY\mid \sG] = Y\E[X\mid\sG]$ a.s.
\item \textbf{Tower property:} If $\sH \subseteq \sG$, then
$\E[\E[X\mid\sG]\mid \sH] = \E[X\mid \sH]$ a.s.
\item \textbf{Irrelevance of independent information:} If $\sH$ is independent of $\sigma (\sG, \sigma (X))$
then
$\E[X|\sigma (\sG, \sH)] = \E[X\mid\sG] \text{ a.s. }$
In particular, if $X$ is independent of $\sH$, then $\E[X\mid \sH] = \E[X]$ a.s.
\item \textbf{$L^2$-projection:} If $\E\abs{X}^2 < \infty$, then $\zeta^\ast \triangleq \E[X\mid\sG]$ minimizes $\E[(X - \zeta)^2]$
over all $\sG$-measurable random variables $\zeta$ such that $\E\abs{\zeta}^2 < \infty$. 
\end{enumerate}
\end{prop}
\end{frame}



\begin{frame}
\begin{Proof}[Proof Continued.]
\begin{enumerate}
\setcounter{enumi}{7}
\item \textbf{$L^2$-projection:} 
We define $L^2(\sG) \triangleq \set{\zeta \text{ a } \sG \text{ measurable random variable}: 
\E\zeta^2 < \infty%,\zeta^{-1}(-\infty, x] \in \sG \text{ for all }x\in \R
}$. 
From the conditional Jensen's inequality applied to convex function $(\cdot)^2:\R \to \R_+$, 
we get that $(\E[X\mid\sG])^2 \le \E[X^2\mid\sG]$ a.s.. 
Since $X\in L^2$, it follows that $X^2\in L^1$ and hence $\E[X\mid\sG] \in L^2$. 
It follows that $\zeta^\ast \triangleq \E[X\mid\sG] \in L^2(\sG)$ from the definition of conditional expectation. 
 
We first show that $X - \zeta^\ast$ is uncorrelated with all $\zeta \in L^2(\sG)$. 
Towards this end, we let $\zeta\in L^2(\sG)$ and observe that 
\EQ{
\E[(X - \zeta^\ast)\zeta] 
= \E[\zeta X] - \E[\zeta \E[X\mid\sG]] 
= \E[\zeta X] - \E[\E[\zeta X\mid \sG]] = 0.
}
\end{enumerate}
\end{Proof}
\end{frame}

\begin{frame}
\begin{Proof}[Proof Continued.]
\begin{enumerate}
\setcounter{enumi}{7}
\item \EQ{
\E[(X - \zeta^\ast)\zeta] 
= \E[\zeta X] - \E[\zeta \E[X\mid\sG]] 
= \E[\zeta X] - \E[\E[\zeta X\mid \sG]] = 0.
}
The above equality follows from the linearity of expectation, 
the $\sG$-measurability of $\zeta$,  
and the definition of conditional expectation. 
Since $\zeta^\ast\in L^2(\sG)$, we have $(\zeta-\zeta^\ast)\in L^2(\sG)$. 
Therefore, $\E[(X - \zeta^\ast)(\zeta-\zeta^\ast)] = 0$.  
 
For any $\zeta\in L^2(\sG)$, 
we can write from the linearity of expectation 
\EQ{
\E(X-\zeta)^2 
= \E(X-\zeta^\ast)^2 + \E(\zeta-\zeta^\ast)^2-2\E(X-\zeta^\ast)(\zeta-\zeta^\ast)
\ge \E(X-\zeta^\ast)^2.
}
\end{enumerate}
\end{Proof}
\end{frame}
\end{document}
